%! Absolute Value Functions \& Transformations — Unit 2, Lesson 2.3 — Slides (Beamer)
\documentclass[aspectratio=169, 11pt]{beamer}
\usepackage{algebra2-beamer}   % provides \forestheader{} and \sectionlabel[color]{}
% Coordinate grid: {xmin}{xmax}{ymin}{ymax}
\newcommand{\lgrid}[4]{%
  \draw[step=1, linegray!40, very thin] (#1,#3) grid (#2,#4);
  \draw[->, thick, charcoal] (#1,0)--(#2,0) node[below right, font=\tiny]{$x$};
  \draw[->, thick, charcoal] (0,#3)--(0,#4) node[left, font=\tiny]{$y$};}

\begin{document}

% TITLE SLIDE (CourseName is not defined in beamer, so the name is literal)
{
\setbeamercolor{background canvas}{bg=forest}
\begin{frame}[plain]
  \vspace{2.8cm}\hspace{0.55cm}%
  \begin{minipage}{0.9\textwidth}
    {\color{goldacc}\bfseries\small LESSON 2.3}\\[0.5cm]
    {\color{white}\bfseries\LARGE Absolute Value Functions \& Transformations}\\[1.2cm]
    {\color{forestmid}\small Algebra 2: Shepherd~$\cdot$~Unit 2: Linear Functions}
  \end{minipage}
\end{frame}
}

% HOOK
\begin{frame}[t, plain]
  \forestheader{A distance graph is a V}
  \sectionlabel{Hook}
  \begin{columns}[T]
    \begin{column}{0.54\textwidth}
      You drive toward a rest stop, reach it at $t=2$ h, and keep going. Distance from it (tens of
      miles) is $d(t)=30\,|t-2|$.
      \begin{itemize}\setlength{\itemsep}{4pt}
        \item $d(0)=60$, \; $d(2)=0$, \; $d(4)=60$
        \item Closest at $t=2$: distance $0$ --- the \textbf{vertex}.
      \end{itemize}
    \end{column}
    \begin{column}{0.44\textwidth}
      \centering
      \begin{tikzpicture}[scale=0.42]
        \lgrid{-1}{5}{-1}{7}
        \draw[very thick, forest] (-0.2,6.6)--(2,0)--(4.2,6.6);
        \fill[charcoal] (2,0) circle (3pt);
      \end{tikzpicture}
    \end{column}
  \end{columns}
  \vspace{2pt}
  \begin{block}{Today}
    A distance is never negative, so its graph bottoms out at one turning point and is symmetric about
    the line through it. Every absolute value function is this V, moved and stretched.
  \end{block}
\end{frame}

% PARENT FUNCTION
\begin{frame}[t, plain]
  \forestheader{The parent $f(x)=|x|$}
  \sectionlabel{Parent}
  \begin{columns}[T]
    \begin{column}{0.52\textwidth}
      \begin{itemize}\setlength{\itemsep}{5pt}
        \item Vertex $(0,0)$; \; axis $x=0$; \; minimum $0$.
        \item Domain: all reals. \; Range: $y\ge 0$.
        \item Decreases on $(-\infty,0]$, increases on $[0,\infty)$.
        \item Really two lines: $y=x$ ($x\ge 0$), $y=-x$ ($x<0$).
      \end{itemize}
    \end{column}
    \begin{column}{0.44\textwidth}
      \centering
      \begin{tikzpicture}[scale=0.5]
        \lgrid{-4}{4}{-1}{5}
        \draw[very thick, forest] (-4,4)--(0,0)--(4,4);
        \fill[charcoal] (0,0) circle (3pt);
      \end{tikzpicture}
    \end{column}
  \end{columns}
\end{frame}

% SHIFTS
\begin{frame}[t, plain]
  \forestheader{Up / down and left / right}
  \sectionlabel[goldacc]{Shifts}
  \begin{columns}[T]
    \begin{column}{0.5\textwidth}
      \textbf{Outside} the bars $\to$ vertical (as expected):
      \[ y=|x|+k \ \text{shifts up/down.} \]
      \textbf{Inside} the bars $\to$ horizontal and \emph{backwards}:
      \[ y=|x-h| \ \text{moves the vertex to } x=h. \]
      Tip: the vertex is where the inside $=0$.
    \end{column}
    \begin{column}{0.46\textwidth}
      \centering
      \begin{tikzpicture}[scale=0.44]
        \lgrid{-2}{6}{-1}{5}
        \draw[very thick, forest!45, dashed] (-1,1)--(0,0)--(4,4);
        \draw[very thick, forest] (-1,4)--(3,0)--(6,3);
        \fill[charcoal] (3,0) circle (3pt);
        \node[charcoal, font=\scriptsize] at (3.4,-0.7){$y=|x-3|$};
      \end{tikzpicture}
    \end{column}
  \end{columns}
\end{frame}

% STRETCH / REFLECT
\begin{frame}[t, plain]
  \forestheader{Stretch, compress, reflect: $y=a|x|$}
  \sectionlabel[goldacc]{The number $a$}
  \begin{columns}[T]
    \begin{column}{0.5\textwidth}
      \begin{itemize}\setlength{\itemsep}{6pt}
        \item $|a|>1$: \textbf{narrower}; \; $0<|a|<1$: \textbf{wider}.
        \item $a<0$: \textbf{reflect} --- the V opens \emph{down}.
        \item $a$ is the slope of the right-hand ray.
      \end{itemize}
    \end{column}
    \begin{column}{0.46\textwidth}
      \centering
      \begin{tikzpicture}[scale=0.42]
        \lgrid{-3}{3}{-4}{5}
        \draw[very thick, forest] (-2.2,4.4)--(0,0)--(2.2,4.4);
        \node[forest, font=\scriptsize] at (2.9,4.4){$2|x|$};
        \draw[very thick, navy] (-3,-3)--(0,0)--(3,-3);
        \node[navy, font=\scriptsize] at (3,-3.6){$-|x|$};
      \end{tikzpicture}
    \end{column}
  \end{columns}
\end{frame}

% VERTEX FORM
\begin{frame}[t, plain]
  \forestheader{Vertex form: $g(x)=a|x-h|+k$}
  \sectionlabel{Putting it together}
  \begin{itemize}\setlength{\itemsep}{5pt}
    \item \textbf{Vertex} $(h,k)$; \; \textbf{axis} $x=h$.
    \item Opens \textbf{up} if $a>0$ ($k$ is a \emph{min}); \; \textbf{down} if $a<0$ ($k$ is a \emph{max}).
    \item $|a|$ sets the width.
  \end{itemize}
  \vspace{2pt}
  \begin{block}{Worked: $g(x)=-2|x-1|+3$}
    Vertex $(1,3)$, axis $x=1$, opens \textbf{down} $\Rightarrow$ \textbf{maximum} $3$; narrow.
    And $g(0)=-2|{-1}|+3=1$ straight from the rule.
  \end{block}
\end{frame}

% ROADMAP
\begin{frame}[t, plain]
  \forestheader{Where Unit 2 goes next}
  \sectionlabel{Connections}
  \textbf{Today:} the parent V $y=|x|$; shifts, stretch, and reflection; vertex form
  $a|x-h|+k$ with its vertex, axis, and min/max.\\[8pt]
  \begin{itemize}\setlength{\itemsep}{4pt}
    \item \textbf{2.4} Piecewise-defined functions --- the V is two lines pasted at the vertex.
    \item \textbf{2.5} Linear regression.
  \end{itemize}
  \vspace{4pt}
  {\footnotesize The solutions of $|x-1|=2$ from Lesson 2.2 are where this V meets the line $y=2$.}
\end{frame}

\end{document}
